\begin{titlepage}
\centering
{\Large Notas didácticas de una entrevista a fondo\par}
\vspace{1.0cm}
{\huge\bfseries La escasez de datos en robótica: simulación, datos sintéticos y la cadena industrial de la inteligencia corporizada}\par
\vspace{0.5cm}
{\Large Zhang Xiaojun conversa con Xie Chen (谢晨), fundador de Guanglun Zhineng (光轮智能)\par}
\vspace{0.35cm}
{\large Elaborado a partir del video público de Zhang Xiaojun Podcast, la transcripción por GPU, la estructura de capítulos y diagramas conceptuales propios\par}
\vspace{0.3cm}
{\large 2025-07-15\par}
\vspace{0.85cm}
\includegraphics[width=0.88\textwidth,height=0.42\textheight,keepaspectratio]{cover.jpg}\par
\vfill
\begin{tcolorbox}[width=0.92\textwidth, colback=black!2!white, colframe=black!55, sharp corners]
\textbf{Título del video}: 109. ¿La robótica enfrenta escasez de datos? Conversación con Xie Chen: simulación y datos sintéticos, la adquisición multimillonaria de Meta y Alexandr Wang\par
\textbf{Canal del video}: Zhang Xiaojun Podcast\par
\textbf{Duración del video}: 01:41:09\par
\textbf{Enlace del video}: \href{\videourl}{\nolinkurl{\videourl}}\par
\textbf{Notas sobre el material}: YouTube no tiene subtítulos disponibles; el texto principal se elaboró a partir de la transcripción con faster-whisper large-v3 en CUDA, los capítulos del video, la descripción del video, la portada y figuras didácticas propias. El video de título fijo no se reutiliza en el texto principal.\par
\textbf{Página del proyecto}: \href{\repourl}{\nolinkurl{\repourl}}\par
\end{tcolorbox}
\end{titlepage}

\tableofcontents
\newpage

